\documentclass[11pt]{article}
\usepackage[a4paper,margin=1in]{geometry}
\usepackage{amsmath,amssymb,amsthm,booktabs,hyperref}
\newtheorem{theorem}{Theorem}
\newtheorem{proposition}{Proposition}
\newtheorem{lemma}{Lemma}
\newtheorem{definition}{Definition}
\title{Supplementary proofs for ``From ecological states to distinguishable futures''}
\author{Author names withheld for drafting}
\date{}

\begin{document}
\maketitle

\section{Scope and notation}
This Supplement proves the four mathematical claims used as the spine of Paper B. The finite latent-world set is denoted by $W$. A current deterministic experiment supplies a record map $R:W\to\mathcal X$. A declared target is $T:W\to\mathcal Y$. A finite action set $A$ acts through deterministic successor maps $F_a:W\to W$. The stochastic reporting extension uses a likelihood $L(x\mid w)$ and the positive-support compatible set
\[
C(x)=\{w\in W:L(x\mid w)>0\}.
\]
These support-level results are distinct from the posterior false-resolution rule used in the stochastic benchmark.

\section{Result 1: experiment-induced information and honest reports}
\begin{definition}[Experiment partition]
For a deterministic record map $R$, define
\[
w\sim_R w' \quad\Longleftrightarrow\quad R(w)=R(w').
\]
Its equivalence classes form the experiment-induced partition $\Pi_R$.
\end{definition}

\begin{proposition}[Record factorization]
Every deterministic function $g:W\to Z$ that is computable from the realized record alone factors through $\Pi_R$. Equivalently, $g$ is constant on every block of $\Pi_R$.
\end{proposition}

\begin{proof}
If $g$ is computed from the record, then there exists $h:\mathcal X\to Z$ such that $g=h\circ R$. If $w\sim_R w'$, then $R(w)=R(w')$, hence
\[
g(w)=h(R(w))=h(R(w'))=g(w').
\]
Thus $g$ is constant on each record class.
\end{proof}

\begin{theorem}[Exact deterministic target-report criterion]
A deterministic record-based report equal to $T(w)$ for every $w\in W$ exists if and only if $T$ is constant on every block of $\Pi_R$.
\end{theorem}

\begin{proof}
Necessity follows from the factorization proposition with $g=T$. For sufficiency, if $T$ is constant on each record class, define $h(x)$ to be that common target value on the class $R^{-1}(x)$. Then $T=h\circ R$.
\end{proof}

\begin{proposition}[Sharp set-valued report]
For realized record $x$ in the image of $R$, the exact set of target values not excluded by the deterministic experiment is
\[
\mathcal T(x)=\{T(w):R(w)=x\}.
\]
No exact record-based procedure can eliminate an element of $\mathcal T(x)$ without an additional assumption or observation.
\end{proposition}

\begin{proof}
Every $y\in\mathcal T(x)$ is attained by at least one world that produces the same record $x$. A report excluding $y$ would therefore be false in that compatible world while using the same observed input. Conversely, every target value outside $\mathcal T(x)$ is absent from all worlds generating $x$ and is excluded by the record itself.
\end{proof}

\begin{proposition}[Stochastic positive-support extension]
For a likelihood model $L(x\mid w)$, a deterministic report guaranteed correct for every world that can produce realized record $x$ exists if and only if $T$ is constant on
\[
C(x)=\{w:L(x\mid w)>0\}.
\]
Otherwise the sharp support-level target report is $T(C(x))$.
\end{proposition}

\begin{proof}
The proof is the same compatibility argument as above with the deterministic preimage $R^{-1}(x)$ replaced by the positive-likelihood support $C(x)$.
\end{proof}

The stochastic benchmark in the main text uses a different contract: after Bayesian updating, a singleton target may be reported when its posterior error is below a declared false-resolution limit. That posterior-risk rule is deliberately not identified with the exact support-level theorem above.

\section{Result 2: unique coarsest target-safe quotient}
Let $\Pi$ be a partition of $W$. Write $[w]_\Pi$ for the block containing $w$.

\begin{definition}[Target-safe partition]
A partition $\Pi$ is target-safe for $(R,T,\{F_a\}_{a\in A})$ when:
\begin{enumerate}
\item if $w,w'$ share a block, then $R(w)=R(w')$;
\item if $w,w'$ share a block, then $T(w)=T(w')$;
\item if $w,w'$ share a block, then $F_a(w)$ and $F_a(w')$ share a block for every $a\in A$.
\end{enumerate}
\end{definition}

Define $\Pi_0=\Pi_R$. Given $\Pi_n$, define $\Pi_{n+1}$ by giving each world the signature
\[
\sigma_n(w)=\left([w]_{\Pi_n},\,T(w),\,([F_a(w)]_{\Pi_n})_{a\in A}\right)
\]
and putting two worlds in the same block exactly when their signatures agree.

\begin{lemma}[Monotone finite refinement]
For every $n$, $\Pi_{n+1}$ refines $\Pi_n$. The sequence reaches a fixed point after finitely many steps.
\end{lemma}

\begin{proof}
The first component of $\sigma_n(w)$ is the current block $[w]_{\Pi_n}$, so worlds in different current blocks cannot merge; hence $\Pi_{n+1}$ refines $\Pi_n$. A strict refinement increases the number of blocks, which is at most $|W|$. Thus only finitely many strict refinements are possible.
\end{proof}

Let $\Pi_*$ denote the first fixed point.

\begin{theorem}[Existence and validity]
$\Pi_*$ is target-safe.
\end{theorem}

\begin{proof}
Because every $\Pi_n$ refines $\Pi_0=\Pi_R$, each fixed-point block lies inside one current-record block. Equality of fixed-point signatures gives equality of targets and equality of successor blocks under every declared action. These are exactly the three target-safety conditions.
\end{proof}

\begin{theorem}[Coarseness and uniqueness]
Every target-safe partition $\Gamma$ refines $\Pi_*$. Therefore $\Pi_*$ is the unique coarsest target-safe partition.
\end{theorem}

\begin{proof}
We prove by induction that every target-safe $\Gamma$ refines every $\Pi_n$. By observation preservation, $\Gamma$ refines $\Pi_0$. Suppose it refines $\Pi_n$, and let $w,w'$ share a $\Gamma$ block. Target safety gives $T(w)=T(w')$. For every action $a$, target safety places $F_a(w)$ and $F_a(w')$ in the same $\Gamma$ block; by the induction hypothesis their successors therefore lie in the same $\Pi_n$ block. Since $w,w'$ also lie in the same $\Pi_n$ block, their full signatures $\sigma_n$ agree. Thus they share a $\Pi_{n+1}$ block. Hence $\Gamma$ refines all $\Pi_n$ and in particular $\Pi_*$. If two coarsest target-safe partitions existed, each would refine the other, so their blocks would coincide.
\end{proof}

\begin{corollary}[Preservation under finite action words]
If $w,w'$ share a block of $\Pi_*$, then for every finite declared action word $a_1\cdots a_\ell$, the corresponding successor worlds share a $\Pi_*$ block and have equal target values.
\end{corollary}

\begin{proof}
Induct on word length. The empty word is immediate. Action stability preserves the shared quotient block after one additional action; target constancy in every quotient block then preserves the target report.
\end{proof}

\subsection{Executable construction}
The repository function \texttt{minimal\_target\_safe\_quotient} implements the fixed-point iteration above. Its oracle test enumerates every set partition of a finite five-world witness and verifies that every valid target-safe partition refines the returned fixed point. This finite exhaustive check is not the proof of the theorem; it is an implementation-level guard that the code matches the proved refinement order and minimality claim.

\section{Result 3: lower-bound failure contracts and worst-case guarantee ceilings}
We formalize the probability direction used in the main text. Consider $k$ truly present coordinates, $m$ independent observation modes, and $r$ reads per coordinate within each mode. Suppose every mode operates with probability at least $a>0$. Conditional on operation, every read has sensitivity at least $p>0$, reads are independent across coordinates and repetitions, and false positives are excluded. A mode failure makes every read assigned to that mode negative.

Let
\[
q=(1-p)^r.
\]
In the least-favourable admissible model, every availability equals $a$ and every sensitivity equals $p$.

\begin{proposition}[Exact least-favourable joint-detection frontier]
Under the least-favourable model, the probability of detecting all $k$ coordinates at least once is
\[
G_{k,m,r}(a,p)=
\sum_{s=0}^{k}(-1)^s \binom{k}{s}
\left(1-a+a q^s\right)^m.
\]
For every model satisfying the declared lower bounds, the realized joint-detection probability is at least $G_{k,m,r}(a,p)$.
\end{proposition}

\begin{proof}
For a subset $S$ of $s$ coordinates, one mode leaves every coordinate in $S$ undetected either when the mode fails, with probability $1-a$, or when it operates and all $rs$ corresponding reads miss, with probability $a q^s$. Independence across modes gives
\[
\Pr\!\left(\bigcap_{i\in S}E_i\right)=
(1-a+a q^s)^m,
\]
where $E_i$ is the event that coordinate $i$ is never detected. Inclusion--exclusion over the events $E_1,\ldots,E_k$ gives the displayed formula.
\end{proof}

\begin{proposition}[Worst-case guarantee ceiling]
For fixed $k,m,a,p$ with $a,p>0$,
\[
\lim_{r\to\infty}G_{k,m,r}(a,p)=1-(1-a)^m.
\]
Consequently no joint-detection guarantee valid uniformly over the lower-bound availability contract can exceed $1-(1-a)^m$ using within-mode repetition alone.
\end{proposition}

\begin{proof}
As $r\to\infty$, $q=(1-p)^r\to0$. In the inclusion--exclusion formula, the $s=0$ term equals one, and every $s\ge1$ term converges to $(-1)^s\binom{k}{s}(1-a)^m$. Therefore
\[
\lim_{r\to\infty}G_{k,m,r}(a,p)
=1+(1-a)^m\sum_{s=1}^{k}(-1)^s\binom{k}{s}
=1-(1-a)^m,
\]
because $\sum_{s=0}^{k}(-1)^s\binom{k}{s}=0$ for $k\ge1$.

The bound is sharp for the declared contract because the admissible least-favourable model with every availability exactly $a$ attains this limit. If the true availabilities exceed $a$, realized detection may exceed this value; hence it is a worst-case guarantee ceiling, not a universal realized-probability ceiling.
\end{proof}

\begin{corollary}[Necessary mode floor]
To certify a uniform joint-detection guarantee at least $c\in(0,1)$ using independent modes each available with probability at least $a\in(0,1)$, it is necessary that
\[
1-(1-a)^m\ge c,
\]
equivalently
\[
m\ge \left\lceil\frac{\log(1-c)}{\log(1-a)}\right\rceil.
\]
\end{corollary}

\section{Result 4: finite risk-limited adaptive design}
Let $\mathcal P$ be a finite declared family of adaptive policies. A policy induces a finite set of terminal histories. Under a declared prior and observation kernels, enumerate each terminal history's probability, accumulated cost, posterior target masses, and terminal report.

A deterministic target report is permitted only when its posterior error probability is at most $\alpha$; otherwise the terminal report remains set-valued. Define for policy $\pi$:
\[
P_{\mathrm{correct}}(\pi),\qquad
P_{\mathrm{wrong}}(\pi),\qquad
P_{\mathrm{ambiguous}}(\pi),\qquad
C(\pi).
\]

\begin{theorem}[Finite least-cost feasible policy]
Fix declared bounds
\[
P_{\mathrm{wrong}}(\pi)\le \alpha,
\qquad
P_{\mathrm{ambiguous}}(\pi)\le \beta,
\qquad
C(\pi)\le B.
\]
If at least one policy in finite family $\mathcal P$ satisfies the bounds, then a least-cost feasible policy exists.
\end{theorem}

\begin{proof}
The feasible subset of the finite set $\mathcal P$ is finite and nonempty. The real-valued cost function therefore attains a minimum on that subset.
\end{proof}

The theorem is intentionally modest: it is a finite existence statement, not a universal sequential-design optimality theorem. The substantive reporting convention is that policies may stop with honest ambiguity rather than force a deterministic target when the declared risk contract is not met.

\section{Implementation correspondence}
The proof objects map directly to repository components:
\begin{center}
\begin{tabular}{ll}
\toprule
Mathematical object & Implementation / verification \\
\midrule
experiment partition & \texttt{ced/experiment\_quotient.py} \\
target-safe fixed point & \texttt{ced/target\_safe\_quotient.py} \\
minimality oracle & \texttt{tests/test\_target\_safe\_quotient.py} \\
mode-diverse guarantee & \texttt{ced/mode\_detection.py} \\
exact factor failures & \texttt{ced/overlapping\_modes.py} \\
risk-limited policies & \texttt{ced/adaptive\_target\_policy.py} and benchmark scripts \\
formal proof contract & \texttt{tests/test\_paper\_b\_supplement\_contract.py} \\
\bottomrule
\end{tabular}
\end{center}

\section{Assumption boundaries}
The exact claims above are conditional on the declared finite world set, record or likelihood model, action grammar, target map, and failure assumptions. In particular:
\begin{itemize}
\item the finite world set may omit relevant mechanisms;
\item positive-support compatibility is not the same object as posterior risk-limited reporting;
\item a lower bound on availability supports a worst-case guarantee, not an upper bound on realized detection;
\item independence of declared modes is an assumption unless established externally;
\item posterior draws or particles provide exact guarantees only for the chosen finite support;
\item the finite policy theorem does not imply global optimality over unrestricted sequential policies.
\end{itemize}

These boundaries are part of the theorem statements rather than implementation caveats.

\end{document}